\againframe<6,9->{augpd}

\frame{
    \frametitle{Heights of Simplices From the VPD}
    \begin{corollary}{}
        The VPD tells us the multiset of vertex (and simplex) heights.
    \end{corollary}
}

\frame{
\frametitle{Coordinates of One Point}
\begin{block}{}
\centering
    \only<1>{\includegraphics[width=.5\textwidth]{topology/pht/example-1v-0d}}%
    \only<2>{\includegraphics[width=.5\textwidth]{topology/pht/example-1v-1d}}%
    \only<3>{\includegraphics[width=.5\textwidth]{topology/pht/example-1v-2d}}%
\end{block}
}

\frame[label=vertalgo]{
\frametitle{Finding Multiple Points}
\begin{block}{}
\centering
    \only<1>{\includegraphics[width=.5\textwidth]{topology/pht/example-complex}}%
    \only<2>{\includegraphics[width=.5\textwidth]{topology/pht/example-dOne}}%
    \only<3>{\includegraphics[width=.5\textwidth]{topology/pht/example-dTwo}}%
    \only<4>{\includegraphics[width=.5\textwidth]{topology/pht/example-dTwo-With}}%
    \only<5->{\includegraphics[width=.5\textwidth]{topology/pht/example-dThreeVertices}}
\end{block}
}


\frame<1->[label=outline]{
    \frametitle{Oracle-Based Reconstruction of $K \subset \R^d$}
    \begin{columns}
        \column{.45\textwidth}
        \pause
        \begin{block}{Vertices}
            \begin{enumerate}
                \item {$d+1$ directions}
                \item {$\Theta(d)$ directions}
            \end{enumerate}
        \end{block}
        \pause
        \begin{block}{Edges}
            \begin{enumerate}
                \item {$n_0^2-n_0$ directions}
                \item {$O(n_1 \log n_0)$ directions}
            \end{enumerate}
        \end{block}
        \pause
        \begin{block}{High Dimensional Simplices}
            \begin{enumerate}
                \item {$O(nn_0 2^k+d)$ directions}
            \end{enumerate}
        \end{block}
        \footnotetext{$k=dim(K)$; $n_i=\#K_i$; $n_*=\#K$}
        \column{.45\textwidth}
        \pause
        \begin{block}{References}
            \scriptsize
            \begin{itemize}
                \item RLB, BTF, RM, SM, DLM DS, AS, JS, LW.
                    Learning Simplicial Complexes from Persistence Diagrams.
                    CCCG 2018.
                \item RLB, BTF, RM, SM, DLM DS, AS, JS,
                    LW.  Reconstructing Embedded Graphs from
                    Persistence Diagrams. CGTA 2020.
                \item BTF, SM, DLM, AS, LW. Efficient Graph Reconstruction
                    and Representation Using Augmented Persistence Diagrams. CCCG 2022.
                \item BTF, SM, DLM, AS, LW. A Faithful Discretization of the
                    Persistent Homology Transform. ArXiv:1912.12759.
                \item TO, AS. Sweeping Orders for Simplicial Complex Reconstruction
                    ArXiv:2501.01901.
            \end{itemize}
        \end{block}
    \end{columns}
    }



\frame{
    \frametitle{\onslide<3->{(Not)} Combinatorially Equivalent Diagrams}

    \centering
    \begin{columns}
        \column{.45\textwidth}
        \begin{block}{Direction $\uparrow$}
            \centering
            {\includegraphics[height=1in]{topology/pht/upV-dirs}}%
            \only<1>{
            $$a_y=b_y < c_y$$
            $$D(\uparrow)=\{ (a_y,\infty), (a_y,c_y) \} $$
            }%
            \only<2->{
                $$ [1]=[a]=[b] < [2]=[c] $$
            $$D(\uparrow)=\{ ([1] ,\infty), ([1] ,[2]) \} $$
            }%
        \end{block}

        \onslide<3->{
        \column{.45\textwidth}
        \begin{block}{Direction \rotatebox{-20}{$\uparrow$}}
            \centering
            {\includegraphics[height=1in]{topology/pht/upV-dirt}}%
            $$ [1]=[a] < [2]=[b] < [3]=[c]$$
            $$D(\rotatebox{-20}{$\uparrow$})=\{ ([1],\infty), ([2],[3]) \} $$
        \end{block}
    }
    \end{columns}

    \footnotetext{
        McCleary, Patel. Edit distance and persistence diagrams over
            lattices. SIAM Journal on Applied Algebra and Geometry,
            6(2):134–155, 2022.
    }
}

\againframe{triangles}

\frame[label=strat]{
    \frametitle{Combinatorial Persistent Homology Transform}

    \centering
    \begin{columns}
        \column{.45\textwidth}
        \begin{block}{Complex}
            \centering
            \only<1>{\includegraphics[height=1.5in]{topology/pht/upV-cplx}}%
            \only<2->{\includegraphics[height=1.5in]{topology/pht/ex-3d-cplx}}%
        \end{block}

        \column{.45\textwidth}
        \begin{block}{Stratification}
            \centering
            \only<1>{\includegraphics[height=1.5in]{topology/pht/ex-v-lines}}%
            \only<2->{\includegraphics[height=1.5in]{topology/pht/ex-3d-strat}}%
        \end{block}
    \end{columns}%

    \footnote{
        BTF, Patel. Combinatorial Persistent Homology Transform. Published in
        FODS.}
}

\frame{
    \frametitle{Finite Representation}

    \centering
    \onslide<1->{\includegraphics[height=1.5in]{topology/pht/ex-v-lines}}%
    \onslide<1->{\includegraphics[height=1.5in]{topology/pht/ex-3d-strat}}%

    \begin{block}{}
        \begin{itemize}
            \item \pause Known stratification \pause (e.g., given vertices)
            \item \pause One PD per strata
        \end{itemize}
    \end{block}
}

\frame<1->[label=counting]{
    \frametitle{Theory Versus Practice}

    \begin{block}{Theory}
        \begin{itemize}
            \item Reconstruction approach: $O(n_*n2^k+d)$ VPD
            \item One PD per stratum: $\Theta(n^{d-1})$ in $\S^{d-1} \subset \R^d$
            \item $\epsilon$-net approach: many more
            \item but ... less are actually needed
        \end{itemize}
    \end{block}
}

\frame{
    \frametitle{``Small'' Samples Work}

    \begin{block}{}
        \begin{itemize}
            \pause
        \item $64$ directions in $\R^3$ and $162$ in $\R^3$ for
            {\color{blue}MDS} with PD (Turner,
                Mukherjee, Boyer; 2013.)
            \pause
        \onslide<3->{\item $32$ directions in $\R^2$ for
            {\color{blue}classification} with
                PD (Hofer et al.; 2017)\footnote{
                Hofer, Kwitt, Niethammer, and Uhl. ``Deep learning with
                topological signatures.'' NeurIPS 30 (2017).}}
            \pause
        \onslide<4->{\item $72$ directions in $\R^2$ for
            {\color{blue}regressoin/prediction} with SECT (Crawford et
                al.; 2019)\footnote{Lorin Crawford, Anthea Monod, Andrew X Chen, Sayan Mukherjee, and Raúl
                Rabadán. 2019. Predicting Clinical Outcomes in Glioblastoma: an Application of
                Topological and Functional Data Analysis. J. Amer. Statist. Assoc. (2019),
                1–12.}}
        \end{itemize}
    \end{block}
}


\frame{
    \frametitle{Too Many / Too Few}

    \begin{columns}
        \column{.45\textwidth}

        \begin{block}{}
            \begin{itemize}
                \item Shapes in $\R^2$
                \item $10$-$70$ vertices
                \item Random, MNIST, MPEG7
            \end{itemize}
        \end{block}

        \column{.45\textwidth}

        \pause
        \begin{block}{$\epsilon$-Sample}
            Over 100k directions needed.
        \end{block}

        \pause
        \begin{block}{Random Subsampling}
            Even with $>16$k directions, we miss more than~10\% of the~$\Theta(n_0^{d-1})$ strata.
        \end{block}
    \end{columns}

    \footnotetext{Experiments in collaboration with David Millman and CompTaG students.}
}
